\section[Организация удаленного выполнения команд через BLE-мост на базе ESP32]
{\\Организация удаленного выполнения команд через BLE-мост \\ на базе ESP32}

\subsection{Цель работы}
Изучить принципы ретрансляции данных в Bluetooth Low Energy (BLE) сетях. Получить практические навыки создания GATT-сервисов и взаимодействия между устройствами. Реализовать архитектуру «Клиент–Шлюз–Сервер», где микроконтроллер ESP32 выступает в роли BLE-моста между двумя хостами, обеспечивая удаленное выполнение команд.

\subsection{Задачи}
\begin{itemize}[nosep]
    \item Изучить архитектуру BLE и принципы работы GATT (сервисы и характеристики).
    \item Рассмотреть механизм передачи данных через BLE-характеристики.
    \item Реализовать промежуточный узел (ESP32), выполняющий ретрансляцию данных.
    \item Настроить взаимодействие между двумя ПК через ESP32.
    \item Разработать серверную часть на Python для выполнения команд.
    \item Разработать клиентский терминал для отправки команд.
    \item Исследовать задержки передачи данных через BLE.
\end{itemize}

\subsection{Теоретический материал}

\subsubsection{Архитектура BLE и GATT}
В BLE взаимодействие устройств строится на основе модели GATT (Generic Attribute Profile), где:
\begin{itemize}
    \item \textbf{Сервис (Service)} — логическая группа характеристик.
    \item \textbf{Характеристика (Characteristic)} — канал передачи данных.
\end{itemize}

В данной работе ESP32 реализует BLE-сервер с двумя характеристиками:
\begin{itemize}
    \item RX — для приема данных от клиента;
    \item TX — для отправки данных клиенту.
\end{itemize}

\textbf{Примечание.} Более подробно архитектура BLE описана в лабораторной работе №5.

\subsubsection{Архитектура «BLE-мост»}
ESP32 выполняет роль промежуточного узла (шлюза):
\begin{enumerate}[nosep]
    \item Принимает команды от клиента через характеристику RX.
    \item Передает команды серверному ПК.
    \item Получает результат выполнения.
    \item Отправляет ответ обратно клиенту через характеристику TX.
\end{enumerate}

Таким образом реализуется прозрачная двунаправленная передача данных.

\subsubsection{Буферизация и ограничения BLE}
BLE передает данные небольшими пакетами (обычно до 20–247 байт в зависимости от MTU). Поэтому:
\begin{itemize}
    \item длинные сообщения необходимо разбивать на части (chunking);
    \item требуется учитывать задержки между пакетами;
    \item возможны потери синхронизации при высокой нагрузке.
\end{itemize}

Общее время отклика можно оценить как:
\[
T_{total} = T_{BLE} + T_{proc} + T_{BLE}
\]

\subsubsection{Удаленное выполнение команд}
На сервере используется Python и модуль \texttt{subprocess}:
\begin{lstlisting}[language=Python]
result = subprocess.run(cmd, shell=True, capture\_output=True, text=True)
\end{lstlisting}

\textbf{Важно:} выполнение произвольных команд небезопасно и допустимо только в лабораторной среде.

\subsection{Порядок выполнения лабораторной работы}

\textbf{Оборудование:}
\begin{itemize}
    \item ПК 1 (Client)
    \item ПК 2 (Server)
    \item ESP32
\end{itemize}

\begin{enumerate}

\item \textbf{Подготовка ПК-Сервера (Executor):}
\begin{itemize}
  \item Используйте готовый Python-скрипт серверной части, приведённый в листинге~\ref{lst:lab4_server}.
    

    \item Запустите скрипт — он будет ожидать команды от ESP32.
\end{itemize}

\item \textbf{Прошивка ESP32 (Bridge):}
\begin{itemize}
    \item Используйте файл \texttt{bridge.ino}.
    \item Загрузите прошивку в ESP32.
    \item ESP32 будет:
    \begin{itemize}
        \item принимать команды через RX;
        \item отправлять ответы через TX;
        \item работать как BLE-посредник.
    \end{itemize}
\end{itemize}

\item \textbf{Настройка ПК-Клиента (Client):}
\begin{itemize}
  \item Используйте Python-скрипт клиента, приведённый в листинге~\ref{lst:lab4_client}.
    \item Он выполняет:
    \begin{itemize}
        \item поиск ESP32;
        \item подключение;
        \item отправку команд;
        \item получение ответов.
    \end{itemize}
\end{itemize}

\item \textbf{Запуск системы:}
\begin{enumerate}[label=\alph*)]
    \item Запустите сервер на ПК 2.
    \item Включите ESP32.
    \item Запустите клиент на ПК 1.
\end{enumerate}

\item \textbf{Тестирование:}
\begin{itemize}
    \item Выполните команды:
    \begin{itemize}
        \item \texttt{echo test}
        \item \texttt{uname -a}
        \item \texttt{ls}
    \end{itemize}
    \item Проверьте обработку ошибок (неверные команды).
    \item Измерьте задержку ответа.
\end{itemize}

\end{enumerate}

\begin{lstlisting}[caption={Серверная часть на ПК 2}, label={lst:lab4_server}]
import asyncio
import subprocess
from bleak import BleakScanner, BleakClient

# ESP32 name and characteristic UUIDs from firmware
DEVICE_NAME = "ESP32-BLE-Bridge"
CHARACTERISTIC_UUID_TX = "6E400003-B5A3-F393-E0A9-E50E24DCCA9E"

global_client = None

# Function triggered when a command is received from the Client
async def command_handler(sender, data):
    command = data.decode('utf-8', errors='ignore').strip()
    print(f"[+] Command received: '{command}'")

    try:
        # Execute the command in the OS terminal
        # shell=True allows running system commands like dir, echo, ipconfig
        process = subprocess.run(
            command,
            shell=True,
            stdout=subprocess.PIPE,
            stderr=subprocess.PIPE,
            text=True,
            timeout=10 # 10-second timeout to prevent hanging
        )

        # Collect output (stdout first, if empty — use stderr)
        output = process.stdout if process.stdout else process.stderr
        if not output:
            output = "[System]: Command executed, no output."

    except subprocess.TimeoutExpired:
        output = "[-] Error: Command timeout exceeded (10 sec)."
    except Exception as e:
        output = f"[-] Execution error: {str(e)}"

    print(f"[...] Sending response to client ({len(output)} bytes)...")

    # BLE transmits data in packets. If response is long, split into chunks (MTU)
    # Using safe default chunk size of 200 bytes
    chunk_size = 200
    output_bytes = output.encode('utf-8', errors='ignore')

    for i in range(0, len(output_bytes), chunk_size):
        chunk = output_bytes[i:i+chunk_size]
        if global_client and global_client.is_connected:
            await global_client.write_gatt_char(CHARACTERISTIC_UUID_TX, chunk)
            await asyncio.sleep(0.05) # Small delay between packets

async def main():
    global global_client
    print(f"[...] Searching for device {DEVICE_NAME}...")
    device = await BleakScanner.find_device_by_name(DEVICE_NAME)

    if not device:
        print(f"[-] Device {DEVICE_NAME} not found. Check if ESP32 is powered on.")
        return

    print(f"[+] ESP32 found: {device.address}. Connecting...")

    async with BleakClient(device) as client:
        global_client = client
        print("[+] Successfully connected to the bridge!")
        print("=== Server is running and waiting for commands ===")

        # Subscribe to TX characteristic (ESP32 forwards client commands here)
        # Since bleak requires a sync function in start_notify, we use asyncio.ensure_future
        await client.start_notify(
            CHARACTERISTIC_UUID_TX,
            lambda sender, data: asyncio.ensure_future(command_handler(sender, data))
        )

        # Keep the script running
        while True:
            await asyncio.sleep(1)

if __name__ == "__main__":
    try:
        asyncio.run(main())
    except KeyboardInterrupt:
        print("\n[-] Server stopped.")
\end{lstlisting}

\begin{lstlisting}[caption={Клиентская часть на ПК 1}, label={lst:lab4_client}]
import asyncio
from bleak import BleakScanner, BleakClient

# ESP32 name and characteristic UUIDs from firmware
DEVICE_NAME = "ESP32-BLE-Bridge"
CHARACTERISTIC_UUID_RX = "6E400002-B5A3-F393-E0A9-E50E24DCCA9E"

# Function to handle incoming responses from the server
def notification_handler(sender, data):
    print(f"\n[Response from Server]:\n{data.decode('utf-8', errors='ignore')}")
    print("\nEnter command: ", end="", flush=True)

async def main():
    print(f"[...] Searching for device {DEVICE_NAME}...")
    device = await BleakScanner.find_device_by_name(DEVICE_NAME)

    if not device:
        print(f"[-] Device {DEVICE_NAME} not found. Check if ESP32 is powered on.")
        return

    print(f"[+] ESP32 found: {device.address}. Connecting...")

    async with BleakClient(device) as client:
        print("[+] Successfully connected to the bridge!")

        # Subscribe to notifications to receive responses from the server
        await client.start_notify(CHARACTERISTIC_UUID_RX, notification_handler)

        print("\n=== Client is ready ===")
        print("Example commands: 'dir', 'echo Hello', 'ipconfig' (or 'ifconfig' for Linux)")
        print("Type 'exit' to quit")

        while True:
            # Read command from console (run in a thread to avoid blocking async loop)
            command = await asyncio.to_thread(input, "Enter command: ")

            if command.strip().lower() == 'exit':
                break

            if not command.strip():
                continue

            # Send command to the server via ESP32
            print(f"[...] Sending: '{command}'")
            await client.write_gatt_char(CHARACTERISTIC_UUID_RX, command.encode('utf-8'))

            # Small delay to allow processing and response output
            await asyncio.sleep(0.5)

if __name__ == "__main__":
    try:
        asyncio.run(main())
    except KeyboardInterrupt:
        print("\n[-] Client stopped.")
\end{lstlisting}

\subsection{Содержание отчета}
\begin{enumerate}
    \item Цель работы
    \item Схема: ПК1 $\leftrightarrow$ ESP32 $\leftrightarrow$ ПК2
    \item Листинги:
    \begin{itemize}
        \item код ESP32 (bridge.ino);
        \item сервер;
        \item клиент;
        \item примеры выполнения команд;
        \item таблица задержек;
    \end{itemize}
    \item Ответы на вопросы
    \item Выводы
\end{enumerate}

\section*{Контрольные вопросы}
\begin{enumerate}
    \item Какие ограничения имеет BLE по размеру пакета?
    \item Почему требуется разбиение данных на чанки?
    \item Какие риски связаны с subprocess?
    \item Что произойдет при разрыве соединения?
    \item В чем отличие BLE от классического Bluetooth (SPP)?
    \item Как влияет MTU на производительность?
\end{enumerate}
